% ===== % Essential Document Setup % ===== %
\documentclass[14pt]{article} % Base document class 
\usepackage[utf8]{inputenc} % Input encoding 
\usepackage[a4paper,top=3cm,bottom=1.5cm,left=2.5cm,right=2.2cm]{geometry}% Page margins 
% \usepackage[11pt]{extsizes} % Font size customization

% ===== % Fonts & Typography % ===== %
\usepackage{anyfontsize}
\AtBeginDocument{\fontsize{12pt}{15pt}\selectfont}
\usepackage{mathpazo} % Palatino font 
\usepackage{lmodern} % Improved version of Computer Modern 
\usepackage{setspace} % Line spacing control

% ===== % Colors and Graphics % ===== %
\usepackage{xcolor} % Color definitions 
\usepackage{graphicx} % Include graphics 
\graphicspath{{figures/}}
\usepackage{grffile} % Extended filename support for graphics 
\usepackage{tikz} % TikZ for drawing

% ===== % Math Packages % ===== %
\usepackage{amsmath, amsfonts, amssymb} % AMS math packages 
\usepackage{siunitx} % SI units formatting
\usepackage{braket}

% ===== % Lists and Itemization % ===== %
\usepackage{enumitem} % Custom list formatting

% ===== % Code Listings % ===== %
\usepackage{minted} % Code highlighting 
\usepackage{listings} % Code formatting 
\usepackage{fancyvrb} % Enhanced verbatim support

% ===== % Table and Figure Tools % ===== % 
\usepackage{caption} % Custom captions 
\usepackage{subcaption} % Sub-figures 
\usepackage{float} % Float placement control 
\usepackage{multirow} % Multiple row cells in tables 
\usepackage{multicol} % Multiple columns 
\usepackage{adjustbox} % Adjust content inside a box 
\usepackage{booktabs} % Professional tables

% ===== % Boxes and Styling % ===== %
\usepackage{framed} % Simple framed boxes
\usepackage{tcolorbox} % Colored and styled boxes 
\tcbuselibrary{skins, breakable} % Enable skin and breakable features

% ===== % Define the page style % ===== %
\usepackage{fancyhdr} % Customizing headers and footers
\usepackage{everypage} % to execute code on every page

\setlength{\headheight}{0pt} % Setting the height of the header
\addtolength{\topmargin}{0.4cm} % Adjusting the top margin
\addtolength{\footskip}{-0.2cm} % Adjusting the top margin

\fancypagestyle{customstyle}{
  \fancyhf{} % Clear header and footer
  \fancyhead[L]{Experiment $A_4$} % Left header
  \fancyhead[R]{Page \thepage\ of 15} % Right header with page number
  \renewcommand{\headrulewidth}{0.4pt} % Horizontal line under the header
}
% Apply the custom page style
\pagestyle{customstyle}

\fancypagestyle{customstyle}{
\fancyhf{} % Clear header and footer
\fancyhead[L]{\textit{Quantum Field Theory}} % Left header
\fancyhead[R]{Exercises}

\fancyfoot[C]{\textbf{\textit{\thepage}}} % Right header with page number
\renewcommand{\headrulewidth}{0.4pt} % Horizontal line under the header
\renewcommand{\footrulewidth}{0.4pt}
}

% ===== % Apply the custom page style % ===== %
\pagestyle{customstyle}

% ===== % Table of Contents & Navigation % ===== % 
\usepackage{hyperref} % Hyperlinks support 
\hypersetup{colorlinks=true, allcolors=black} % Black hyperlinks 
\usepackage{tocloft} % Custom TOC formatting 
\renewcommand{\cftsecleader}{\cftdotfill{\cftdotsep}} % Dotted TOC lines 
\usepackage{titletoc} % TOC appearance customization

% ===== % Citations % ===== % 
\usepackage[numbers]{natbib} % Numbered citations

% ===== % Miscellaneous % ===== %
\usepackage{etoolbox} % Logic and patching commands 
\usepackage{changepage} % Adjustable width environment 
\usepackage{lipsum} % Dummy text

\begin{document}
\begin{titlepage}
    \centering
    \vspace*{3cm}
    {Exercises on \par}
    \vspace{1cm}
    \hrule
    \vspace{1cm}
    {\huge \textit{\textbf{Quantum Field Theory}} \par}
    \vspace{1cm}
    \hrule
    \vspace{1cm}
    \small
   (Based on the lectures by \textit{Dr. M Arshad Momen})
\vspace{3cm}
% \vfill
\end{titlepage}
%-------------------------------
% Table of Contents
%-------------------------------
\tableofcontents
\newpage
\section{Exercise for Class-2 [Lecture--1]}
Derive the Euler–Lagrange equations corresponding to this action:
\[
S[x] = m \int \sqrt{ \eta_{ab} \, dx^a dx^b }=m \int d \tau \\
    = m \int \sqrt{ \eta_{ab} \, \frac{dx^a}{d\lambda} \frac{dx^b}{d\lambda} } \, d\lambda
\]
\section{Exercise for Class-4 [Lecture-1]}
\subsection{Assignment-1}
How does $\hat{J}^{rs}$ transform under Poincaré transformation? (It's a semi direct product)

\subsection{Assignment-2}
Recover $3D$ rotation algebra from angular momentum.

\section{Exercise for Class-5 [Lecture-2]}
\subsection{Exercise-1}
Show that we can recover the Lorentz force from the action:
\[
S = \frac{1}{2} \int m~ \dot{\xi}^2 d\tau + q \int A_b(\xi) \dot{\xi}^b \, d\tau
\]
\[
\left[~\text{This latter interaction term is geometrical as well}~\right]
\]

\subsection{Exercise-2}
Check the commutations:
\begin{align*}
    [J, p]~~ &\sim ~~p \\
    [p, p]~~ &\sim ~~0 \\
    [J, J]~~ &\sim ~J \\
    [J, p^2]~ &\sim ~0 \\
    [J, p_ap^a] &= [J, p_a]p^a + p_a[J, p^a]\sim 0
\end{align*}



\section{Exercise for Class-6 [Lecture-3]}
\subsection{Exercise-1}
Show that the lie derivatives of vector field satisfies:
\[
\mathfrak{L}_{\hat{V}} \mathfrak{L}_{\hat{W}} - \mathfrak{L}_{\hat{W}} \mathfrak{L}_{\hat{V}} = \mathfrak{L}_{[\hat{V},\hat{W}]}
\]
\subsection{Exercise-2}
Derive the lie derivative for a tensor $\mathfrak{L}_{\hat{V}}\, \tilde{T}^a\,_b $.

\section{Exercise for Class-7 [Lecture-4]}
$\bullet$ Calculate $\mathfrak{L}~g_{ab}$ \\

\textbf{Hint :} Let choose, 
\[ g_{ab} = P_a Q_b \]
\begin{align*}
& \mathfrak{L}_v (P_a Q_b) = (\mathfrak{L}_vP_a) Q_b + P_a (\mathfrak{L}_v Q_b) \\
& \mathfrak{L}_v g_{ab} = V^c~ \partial_c\, g_{ab} + g_{ac}\, \partial_b V^c + g_{cb}\, \partial_a V^c
\end{align*}
$\bullet$ Verify if the partial derivative is switched with the covariant derivative ($\partial \rightarrow\nabla$) the Lie derivative won't change.

\section{Exercise for Class-8 [Lecture-5]}
Derive the equation of motion for this action and compare with the KG equation to determine the constants -
\[
S[\phi] = \alpha\, \phi_i \phi_i + \beta\, \partial \phi_i \, \partial \phi_i + \gamma_i\, \phi_i + K_i \partial \phi_i
\]
\textbf{Note:} The last term is important for $\pi$-meson.
\section{Exercise for Class-11 [Lecture-8]}
Find the solution of the G(t).
\section{Exercise for Class-13 [Lecture-10]} 
Do this calculation for 6-point function and calculate $\langle \phi_a \phi_b \phi_c \phi_d \phi_e \phi_f\rangle $.
\section{Exercise for Class-14 [Lecture-11]} 
\subsection{Exercise-1}
Consider a field theory for two fields $\phi, \chi$ and whose action is given by
\begin{align*}
S = \frac{1}{2} \phi^T A \phi + \frac{1}{2} \chi^T A\, \chi + \chi^T B \, \phi - \phi^T B\, \chi \quad, \quad [B = -B]
\end{align*}
Find the propagators for this field.
\subsection{Exercise-2}
For $V=\frac{{g \phi^3}}{3!}$, Calculate the $2^{\text{nd}}$ order contribution to $Z[J]$ due to this interaction.
\section{Exercise for Class-15 [Lecture-12]} 
Decompose the full 4-point function in terms of the connected Green's function of 4th order or less.
\section{Exercise for Class-16 [Lecture-13]} 
Decompose the 4-point connected Green function in terms of proper vertices and the 2-point function.
\section{Exercise for Class-19 [Lecture-16]} 
\begin{framed}
For $n$ such vertex insertions: each minus comes from the propagator 
in a Euclideanized way,
\[
(-1)^n \int \frac{d^Dk}{(2\pi)^D} \frac{[V''(\varphi)]^n}{(k^2+m^2)^n}.
\]

Using dimensional regularization,
\[
\int \frac{d^Dk}{(2\pi)^D} \frac{1}{(k^2+m^2)^n} 
= \frac{1}{(4\pi)^{D/2}} \frac{\Gamma\!\left(n-\tfrac{D}{2}\right)}{\Gamma(n)} 
\frac{1}{(m^2)^{\,n-D/2}}.
\]

Thus,
\[
=(-1)^n \frac{[V''(\varphi)]^n}{(4\pi)^{D/2}}
\frac{1}{\Gamma(n)} \frac{\Gamma\!\left(n-\tfrac{D}{2}\right)}{(m^2)^{\,n-D/2}}.
\]

So,
\[
\Gamma[\varphi] = \frac{\hbar}{2i} \sum_{n=1}^{\infty} 
\left[ -\frac{V''(\varphi)}{m^2} \right]^n 
\frac{\Gamma\!\left(n-\tfrac{D}{2}\right)}{\Gamma(n)} 
(m^2)^{D/2}.
\]
\end{framed}
\textbf{HW:} Resume this series (the first terms) to obtain a closed expression.
\section{Exercise for Class-20 [Lecture-17]} 
\subsection{Assignment-1}
Redo today's calculation for $D=3,6$.
\subsection{Assignment-2}
    Fix the today's $\mu$ factor.
\subsection{Assignment-3}
    Express $V_T$ in terms of $m^-2$ and $\phi$.
    \subsection{Assignment-4}
    \begin{framed}
        \[
\frac{d \Gamma_1}{d m^2}
= \frac{\hbar}{16 \pi^2}
\int_0^{\Lambda^2} \frac{z}{z + m^2} dz
= \frac{\hbar}{16 \pi^2}
\int_0^{\Lambda^2} \left[ 1 - \frac{m^2}{z + m^2} \right] dz.
\]
\end{framed}
Perform the integral explicitly.
\subsection{Assignment-5}
Check that 
\[
[Q^a, Q^b] = i f^{abc} Q^c
\]
\subsection{Assignment-6}
\[
\int d\eta_1 \dots d\eta_N \; e^{\eta^T A \eta} = \det(A)
\]

By induction method:
\[
\eta^T = -\eta, \quad A^T = -A
\]
\end{document}
